% Part: second-order-logic
% Chapter: metatheory
% Section: loewenheim-skolem

\documentclass[../../../include/open-logic-section]{subfiles}

\begin{document}

\olfileid[fr]{sol}{met}{lst}

\olsection{Les théorèmes de Löwenheim--Skolem
  échouent au second ordre}

\begin{explain}
Pour un langage du premier ordre au plus
dénombrable, le théorème de Löwenheim--Skolem
\emph{descendant} affirme que tout ensemble
d'!!{sentence}s ayant un modèle infini possède
un modèle !!{denumerable}. Sa démonstration
peut s'appuyer sur le théorème de complétude :
en ajoutant une infinité dénombrable de constantes contraintes à être
distinctes, on obtient un modèle infini au plus
dénombrable. Il existe aussi un théorème de
Löwenheim--Skolem \emph{ascendant} : un ensemble
d'!!{sentence}s qui possède un modèle
!!{denumerable} possède aussi un modèle
!!{nonenumerable}. Sous la sémantique standard,
ces deux théorèmes échouent pour la logique du
second ordre.
\end{explain}

\begin{thm}
\ollabel{thm:sol-no-ls}
Le théorème de Löwenheim--Skolem descendant
échoue au second ordre : certains !!{sentence}s
ont des modèles infinis, mais aucun modèle
!!{enumerable}.
\end{thm}

\begin{proof}
Rappelons que
\[
\fn{Count} \ident \lexists[z][\lexists[u][\lforall[X][((X(z) \land
      \lforall[x][(X(x) \lif X(u(x)))]) \lif \lforall[x][X(x)])]]]
\]
est vrai dans une !!{structure}~$\Struct{M}$
si et seulement si $\Domain{M}$ est
!!{enumerable}. Par conséquent,
$\lnot\fn{Count}$ est vrai dans~$\Struct{M}$
si et seulement si ce !!{domain} est
!!{nonenumerable}. De tels !!{domain}s
existent : prenons, par exemple,
$\Pow{\Nat}$ ou~$\Real$.
Ainsi, $\lnot\fn{Count}$ possède des modèles
infinis, mais aucun modèle !!{enumerable}.
\end{proof}

\begin{thm}
Il existe des !!{sentence}s ayant des modèles
!!{denumerable}s, mais aucun modèle
!!{nonenumerable}.
\end{thm}

\begin{proof}
$\fn{Count}\land\fn{Inf}$ est vrai dans
$\Nat$, mais dans aucune !!{structure}
dont le !!{domain} est !!{nonenumerable}.
\end{proof}

Le raisonnement descendant utilise le fait que le langage enrichi
reste au plus dénombrable. Dans le sens ascendant, on peut
au contraire choisir une famille non dénombrable de nouvelles
constantes et imposer qu'elles soient deux à deux distinctes.
Chaque sous-ensemble fini des contraintes est satisfaisable
dans un modèle infini de la théorie. Par compacité, il existe
un modèle de l'ensemble des contraintes, dont le réduit
au langage initial est encore non dénombrable.

\end{document}
